% =============================================================================
% Section 4.6 — Controlled Validation Under Fixed Weather
% Drop-in LaTeX file for Chapter 4 of the thesis.
% =============================================================================

\subsection{Controlled Validation Under Fixed Weather}
\label{sec:results:controlled-validation}

The preceding sections established that charging scheduling policy,
fleet survivability, and network quality form a coherent cross-layer
mechanism under the stochastic weather conditions of Test Round~2.
A natural question is whether the observed trends are primarily driven
by weather variance itself or whether they reflect a more fundamental
property of the scheduling--survivability--network chain.
To answer this, Test Round~3 repeats the protocol comparison under
\emph{fixed sunny weather}, eliminating the Markov regime transitions
that introduced time-varying routing stress in Round~2.
If the cross-layer relationships persist despite the removal of weather
variance, the mechanism can be considered robust to the environmental
factor that dominated Round~2 degradation episodes.

% -----------------------------------------------------------------------------
\subsubsection{Rationale and Experimental Setup}
\label{sec:results:controlled-validation:setup}

\paragraph{Design changes from Round~2 to Round~3.}
The only deliberate change between the two rounds is the weather
configuration.  In Round~2, the weather node operated in Markov mode,
cycling through sunny, cloudy, rainy, and stormy regimes according to
stochastic transition probabilities.  In Round~3, the weather node is
locked to the \texttt{sunny} regime for the entire mission duration.
All \texttt{weather\_timeseries.csv} files in the Round~3 dataset
confirm a single regime label (\texttt{sunny}) throughout every run.
By fixing weather, WEATHER\_DROP-induced routing failures are removed,
which should reduce both ROUTING\_DROP rates and PDR variance across
all protocols.

\paragraph{What remains constant.}
The protocol set (seven policies), fleet composition (two cluster heads,
three member UAVs, one UGV with charging dock), simulated mission
duration ($\sim$180~min per run), number of replicates (three per
protocol), logging schema, and metric definitions are identical to
Round~2.  The analysis pipeline is also the same, ensuring that any
observed difference is attributable to weather control rather than
methodological change.

\paragraph{Dataset scope and filtering.}
The Round~3 dataset contains 25~run folders in total.  Four folders
whose names contain \texttt{stormy} (case-insensitive) are excluded
because they belong to a separate sensitivity check and violate the
fixed-sunny design.  After filtering, 21~sunny runs remain:
7~protocols $\times$ 3~replicates each, matching the Round~2 structure
exactly.

\paragraph{Threats to validity.}
Even with fixed weather, residual stochasticity remains from random
seeds governing packet generation timing, UAV mobility jitter, and
charging-request arrival patterns.  The validation therefore targets
\emph{trend consistency} across rounds---whether the same directional
relationships hold---rather than exact numerical reproduction of
Round~2 values.

% -----------------------------------------------------------------------------
\subsubsection{Results and Comparison with Test Round~2}
\label{sec:results:controlled-validation:results}

This section compares the Round~2 and Round~3 outcomes category by
category, using a structured \emph{trend match / mismatch} framework.
For each category, the Round~2 trend is stated first, followed by
the Round~3 observation, the match verdict, and the supporting figure
references.  A consolidated evidence map is provided at the end of
this subsection (Table~\ref{tab:evidence-map}).

% ---- 1. Network Reliability ------------------------------------------------
\paragraph{1) Network reliability: overall PDR and PDR over time.}

\emph{Round~2 trend.}
\texttt{ugv\_p\_edf} achieves the highest mean PDR (0.630), followed by
\texttt{ugv\_fcfs} (0.616) and \texttt{ugv\_dynamic} (0.595).  The
role-based protocols form the weakest tail, with
\texttt{ugv\_p\_role\_priority} at 0.554.  PDR time-series show high
between-run variance, particularly for \texttt{ugv\_edf} and
\texttt{ugv\_p\_role\_priority}, with episodic drops linked to weather
regime transitions.
(Figures: \texttt{01\_mean\_pdr\_merged.png},
\texttt{01\_network\_pdr\_over\_time\_merged.png},
\texttt{CL\_A\_protocol\_kpi\_overview.png} in
\texttt{test\_round2/figures/}.)

\emph{Round~3 observation.}
All protocols show an absolute PDR increase of 0.07--0.20, consistent
with the removal of weather-induced packet loss.
\texttt{ugv\_p\_role\_priority} rises to rank~1 (0.754),
\texttt{ugv\_p\_edf} drops to rank~2 (0.732), and
\texttt{ugv\_dynamic} falls to rank~7 (0.667).
The Kendall rank correlation between Round~2 and Round~3 PDR rankings
is $\tau = -0.24$, indicating no preservation of the ordinal ranking.
(Figures: \texttt{01\_mean\_pdr\_merged.png},
\texttt{01\_network\_pdr\_over\_time\_merged.png},
\texttt{CL\_A\_protocol\_kpi\_overview.png} in
\texttt{test\_round3/figures/}.)

\emph{Trend match:} \textbf{NO} for protocol ranking.
However, the positive association between charging success and PDR
is preserved at a qualitative level: the protocol with the highest
Round~3 charge success rate (\texttt{ugv\_p\_role\_priority}, 0.598)
is also the Round~3 PDR leader.

% ---- 2. Charging Outcomes ---------------------------------------------------
\paragraph{2) Charging outcomes: success / timeout / routing-drop
composition.}

\emph{Round~2 trend.}
\texttt{ugv\_p\_edf} leads in charge success rate (48.6\%), while
\texttt{ugv\_role\_priority} trails at 38.4\%.  ROUTING\_DROP rates
range from 0.218 (\texttt{ugv\_p\_edf}) to 0.346 (\texttt{ugv\_edf}),
and are a significant negative predictor of PDR ($r = -0.53$).
(Figures: \texttt{02\_charge\_success\_rate\_merged.png},
\texttt{03\_charge\_outcome\_breakdown\_merged.png},
\texttt{CL\_G\_routing\_drop\_timeseries.png} in
\texttt{test\_round2/figures/}.)

\emph{Round~3 observation.}
All protocols improve in success rate under sunny weather, with
\texttt{ugv\_p\_role\_priority} reaching the top (0.598) and
\texttt{ugv\_edf} recovering strongly to rank~2 (0.583).
ROUTING\_DROP rates collapse dramatically to 0.010--0.020 across all
protocols, confirming that WEATHER\_DROP was the dominant cause of
request loss in Round~2.
Timeout rates, however, increase substantially for all protocols
(e.g.\ \texttt{ugv\_dynamic} from 0.230 to 0.434), suggesting that
with longer-surviving fleets and fewer routing drops, queue contention
replaces routing failure as the primary bottleneck.
Kendall $\tau$ for charge success rate ranking: 0.33 (NO match).
(Figures: \texttt{02\_charge\_success\_rate\_merged.png},
\texttt{03\_charge\_outcome\_breakdown\_merged.png},
\texttt{CL\_G\_routing\_drop\_timeseries.png} in
\texttt{test\_round3/figures/}.)

\emph{Trend match:} \textbf{NO} for ranking, but the \emph{structural
finding} that routing drops are weather-coupled is strongly confirmed
by their near-elimination in Round~3.

% ---- 3. Survivability -------------------------------------------------------
\paragraph{3) Survivability: battery depletions.}

\emph{Round~2 trend.}
\texttt{ugv\_p\_edf} has the fewest depletions (20.0 per run),
\texttt{ugv\_edf} the most (45.0).  Higher depletion burden is
negatively associated with PDR ($r = -0.57$).
(Figures: \texttt{03\_dead\_uav\_cumulative\_merged.png},
\texttt{CL\_F\_pdr\_vs\_depletions\_timeseries.png} in
\texttt{test\_round2/figures/}.)

\emph{Round~3 observation.}
Depletions decrease for all protocols, with the most dramatic
reductions in the protocols that suffered worst in Round~2:
\texttt{ugv\_edf} drops from 45.0 to 16.7 (rank~1 in Round~3),
and \texttt{ugv\_p\_role\_priority} from 26.7 to 17.3 (rank~2).
The Kendall $\tau$ is $-0.05$ (NO ranking match).
(Figures: \texttt{03\_dead\_uav\_cumulative\_merged.png},
\texttt{CL\_F\_pdr\_vs\_depletions\_timeseries.png} in
\texttt{test\_round3/figures/}.)

\emph{Trend match:} \textbf{NO} for ranking.  The convergence of
depletion counts across protocols is consistent with the removal of
weather-driven cascade failures that disproportionately affected weaker
protocols in Round~2.  With routing drops near zero, even protocols
that previously lost many requests can now successfully charge,
compressing the survivability gap.

% ---- 4. Cross-Layer Associations --------------------------------------------
\paragraph{4) Cross-layer associations.}

\emph{4a) Success rate vs.\ PDR.}
In Round~2, charge success rate shows a positive association with PDR
(Pearson $r = 0.58$, $p = 0.006$).  In Round~3, the protocol with
the highest success rate (\texttt{ugv\_p\_role\_priority}) is also the
PDR leader.  However, because protocol rankings have shifted,
statistical confirmation requires recomputation of the per-run
correlation on the Round~3 dataset.
(Figures: \texttt{CL\_B\_charging\_vs\_pdr\_scatter.png} in both
rounds; \texttt{CL\_E\_correlation\_bar\_chart.png}.)

\begin{quote}\textbf{TODO:} Recompute Pearson and Spearman correlation
between charge success rate and mean PDR for Round~3 runs ($n = 21$).
Verify whether the positive association remains significant.  If the
correlation is computed in
\texttt{CL\_E\_correlation\_bar\_chart.png} for Round~3, read the
values from that figure and report them here.\end{quote}

\emph{4b) Routing-drop rate vs.\ PDR.}
In Round~2, ROUTING\_DROP rate is a strong negative predictor of PDR
($r = -0.53$).  In Round~3, routing drops collapse to near-zero for
all protocols, so the predictor effectively loses variance.  This
result does not contradict the Round~2 finding; rather, it confirms
that once routing drops are removed (by fixing weather), their
explanatory power over PDR disappears because the mechanism is no
longer active.
(Figures: \texttt{CL\_G\_routing\_drop\_timeseries.png},
\texttt{CL\_E\_correlation\_bar\_chart.png} in both rounds.)

\emph{4c) Depletions vs.\ PDR.}
In Round~2, depletions are negatively associated with PDR
($r = -0.57$).  In Round~3, the compressed depletion range makes the
statistical test less powerful, but the qualitative pattern persists:
protocols with fewer depletions (\texttt{ugv\_edf}, 16.7;
\texttt{ugv\_p\_role\_priority}, 17.3) achieve above-average PDR.
(Figures: \texttt{CL\_F\_pdr\_vs\_depletions\_timeseries.png},
\texttt{CL\_D\_mechanism\_chain.png} in both rounds.)

\emph{Trend match for cross-layer associations:} \textbf{PARTIAL.}
The directional relationships are qualitatively preserved, but the
compressed variance in Round~3 (due to weather control) weakens the
statistical power of the correlations.  The routing-drop pathway is
confirmed by its near-elimination.

% ---- 5. Delay Discussion ----------------------------------------------------
\paragraph{5) Delay discussion (cautious).}

End-to-end delay is computed only on \emph{delivered} packets.  As
discussed in Section~4.5, this introduces a conditioning bias: when
PDR drops, the surviving packet set traverses a sparser, shorter
topology, which can paradoxically \emph{reduce} measured delay.  This
caveat remains relevant in Round~3, although the bias should be
attenuated because PDR is uniformly higher and more stable under
fixed sunny weather.

In Round~2, mean E2E delay ranges from 52.8\,ms
(\texttt{ugv\_p\_dynamic\_score}) to 71.4\,ms
(\texttt{ugv\_p\_role\_priority}).  In Round~3, the range narrows to
48.1--61.6\,ms, with \texttt{ugv\_p\_edf} achieving the lowest delay
(48.1\,ms).  The Kendall $\tau$ for delay ranking is $-0.33$ (NO
match), and the ranking reversal---\texttt{ugv\_p\_role\_priority}
improves from worst to second-best---is consistent with the
conditioning bias: in Round~2, \texttt{ugv\_p\_role\_priority} had
the lowest PDR, so its delivered packets were biased toward short
paths; in Round~3, with much higher PDR, its delay rises less because
the full multi-hop topology is preserved.
(Figures: \texttt{06\_mean\_e2e\_delay\_merged.png},
\texttt{06\_pdr\_vs\_e2e\_scatter.png},
\texttt{CL\_C\_e2e\_delay\_conditioning\_analysis.png},
\texttt{CL\_K\_delay\_robustness\_panel.png} in both rounds.)

\emph{Trend match:} \textbf{NO} for ranking.
This mismatch is expected and \emph{does not undermine} the main
findings.  Delay remains a secondary metric in the evaluation
framework, and its Round~3 behaviour is consistent with the
conditioning-bias explanation from Section~4.5.

% ---- Best and Worst Protocol Check ------------------------------------------
\paragraph{Best and weakest protocol stability.}

\emph{Does the Round~2 best protocol remain best in Round~3?}
\textbf{No.}  \texttt{ugv\_p\_edf}, the Round~2 leader in PDR (0.630),
charge success (48.6\%), and lowest depletions (20.0), drops to rank~2
in PDR (0.732) and rank~3 in success rate (0.576) in Round~3.  It is
overtaken by \texttt{ugv\_p\_role\_priority}, which benefits most from
the removal of weather-driven routing drops.

\emph{Does the Round~2 weak tail remain weak?}
\textbf{Partially.}  \texttt{ugv\_role\_priority} remains the weakest
in charge success rate (0.449) and has the highest depletion count
(25.0) in Round~3, preserving its position at the bottom of the
survivability ranking.  However, its PDR improves enough (0.722) to
move out of last place; \texttt{ugv\_dynamic} (0.667) takes over as
the lowest-PDR protocol.

\emph{Plausible explanation for the reversal of
\texttt{ugv\_p\_role\_priority}.}
In Round~2, the CH-priority paradox (Section~4.4.4) described how
role-based protocols suffered from simultaneous CH withdrawal, which
compounded weather-driven member starvation.  In Round~3, the removal
of weather drops means that members retain connectivity to the UGV
scheduler even when both CHs charge concurrently.  This neutralises
the worst aspect of the CH-synchronisation failure mode, allowing
the preemptive role-priority policy to translate its aggressive CH
servicing into an advantage rather than a liability.  The finding
therefore does not invalidate the CH-priority paradox; instead, it
shows that the paradox is \emph{amplified by weather stress} and can
be mitigated when network conditions are favourable.

% ---- Energy per Session: the one partial match ------------------------------
\paragraph{Energy per session: a stable metric.}

Energy recovered per charging session is the only metric with a
\textbf{PARTIAL} match (Kendall $\tau = 0.81$).
\texttt{ugv\_role\_priority} leads in both rounds ($\sim$65--64\,Wh),
and the overall ranking is largely preserved with only minor
mid-tier reordering.  This stability is expected: per-session energy
depends primarily on dock hardware and session duration, neither of
which changes between rounds.
(Figures: \texttt{02\_energy\_recovered\_merged.png},
\texttt{04\_policy\_radar\_merged.png} in both rounds.)

% ---- Evidence Map -----------------------------------------------------------
\begin{table}[htbp]
\centering
\caption{Evidence map: claims, supporting figures, metric definitions,
and caveats for Section~4.6.}
\label{tab:evidence-map}
\footnotesize
\begin{tabular}{p{3.2cm} p{2.8cm} p{2.8cm} p{2.5cm} p{2.7cm}}
\hline
\textbf{Claim} & \textbf{Round~2 figure(s)} & \textbf{Round~3 figure(s)} & \textbf{Metric} & \textbf{Caveat / Note} \\
\hline

PDR ranking reversal
& \texttt{01\_mean\_pdr\_merged.png}
& \texttt{01\_mean\_pdr\_merged.png}
& Mean PDR (delivered/generated, full run)
& Kendall $\tau = -0.24$; absolute PDR higher in R3 \\

PDR temporal stability improves
& \texttt{01\_network\_pdr\_over\_time\_merged.png}
& \texttt{01\_network\_pdr\_over\_time\_merged.png}
& Windowed PDR time-series
& Reduced variance expected under fixed weather \\

Routing drops collapse in R3
& \texttt{CL\_G\_routing\_drop\_timeseries.png}
& \texttt{CL\_G\_routing\_drop\_timeseries.png}
& Charge-request ROUTING\_DROP rate
& Confirms WEATHER\_DROP was dominant cause \\

Charge success ranking shifts
& \texttt{02\_charge\_success\_rate\_merged.png}; \texttt{03\_charge\_outcome\_breakdown\_merged.png}
& \texttt{02\_charge\_success\_rate\_merged.png}; \texttt{03\_charge\_outcome\_breakdown\_merged.png}
& Charge success rate (\#STARTED / \#requests)
& $\tau = 0.33$; timeouts rise as queue contention replaces routing loss \\

Depletion gap compresses
& \texttt{03\_dead\_uav\_cumulative\_merged.png}; \texttt{CL\_F\_pdr\_vs\_depletions\_timeseries.png}
& \texttt{03\_dead\_uav\_cumulative\_merged.png}; \texttt{CL\_F\_pdr\_vs\_depletions\_timeseries.png}
& Total battery depletions per run
& $\tau = -0.05$; weather removal reduces cascade failures \\

Success--PDR association (qualitative)
& \texttt{CL\_B\_charging\_vs\_pdr\_scatter.png}; \texttt{CL\_E\_correlation\_bar\_chart.png}
& \texttt{CL\_B\_charging\_vs\_pdr\_scatter.png}; \texttt{CL\_E\_correlation\_bar\_chart.png}
& Pearson $r$ / Spearman $\rho$ (success vs.\ PDR)
& R3 per-run correlation not yet recomputed (see TODO) \\

Depletion--PDR association persists
& \texttt{CL\_F\_pdr\_vs\_depletions\_timeseries.png}; \texttt{CL\_D\_mechanism\_chain.png}
& \texttt{CL\_F\_pdr\_vs\_depletions\_timeseries.png}; \texttt{CL\_D\_mechanism\_chain.png}
& Depletions vs.\ PDR
& Compressed range weakens stat.\ power \\

Delay ranking reversal
& \texttt{06\_mean\_e2e\_delay\_merged.png}; \texttt{CL\_C\_e2e\_delay\_conditioning\_analysis.png}
& \texttt{06\_mean\_e2e\_delay\_merged.png}; \texttt{CL\_C\_e2e\_delay\_conditioning\_analysis.png}
& Mean E2E delay (delivered packets only)
& Conditioning bias; $\tau = -0.33$; delay is secondary metric \\

Energy/session ranking stable
& \texttt{02\_energy\_recovered\_merged.png}; \texttt{04\_policy\_radar\_merged.png}
& \texttt{02\_energy\_recovered\_merged.png}; \texttt{04\_policy\_radar\_merged.png}
& Mean Wh per charge session
& $\tau = 0.81$ (PARTIAL match); hardware-dominated \\

CH-priority paradox attenuated
& \texttt{CL\_M\_role\_scheduling\_audit.png}; \texttt{CL\_N\_ch\_sync\_cascade.png}
& \texttt{CL\_M\_role\_scheduling\_audit.png}; \texttt{CL\_N\_ch\_sync\_cascade.png}
& Role-stratified success rate, CH gap
& R3 members retain UGV access even during dual-CH charging \\

\hline
\end{tabular}
\end{table}

% -----------------------------------------------------------------------------
\subsubsection{Interpretation}
\label{sec:results:controlled-validation:interpretation}

\paragraph{What the validation strengthens.}
The controlled round confirms two structural properties of the
cross-layer mechanism.
First, the near-elimination of ROUTING\_DROP in Round~3 (rates of
0.010--0.020 versus 0.218--0.346 in Round~2) validates that
weather-driven routing failure was the dominant cause of request loss
in Round~2.  This confirms the feedback loop proposed in
Section~4.4.2: degraded weather $\to$ WEATHER\_DROP $\to$
ROUTING\_DROP of charge requests $\to$ missed charging $\to$ UAV
depletion $\to$ further network degradation.  When the weather trigger
is removed, the loop is broken at its entry point, and all protocols
improve.
Second, the qualitative direction of the charging--survivability--PDR
association persists: protocols that achieve higher charge success
rates tend to show fewer depletions and higher PDR, even though the
specific protocol rankings differ.  The cross-layer mechanism is
therefore not an artefact of weather variance.

\paragraph{What the validation does not prove.}
This section does not establish strict causal claims.  Both Round~2
and Round~3 are observational simulation studies, not controlled
experiments with randomised protocol assignment at the run level.
Although the design holds all factors constant except the scheduling
policy (and, across rounds, the weather mode), confounding from
residual stochastic factors cannot be fully excluded.
The fleet size (five UAVs, one UGV) is also small, limiting
generalisability to larger deployments.  Finally, the simulator
abstracts several physical-layer effects (e.g.\ exact RF propagation,
wind-induced UAV drift), so the magnitudes reported here should not
be taken as field-deployable predictions.

\paragraph{Explaining the mismatches.}
The widespread ranking reversals (9 out of 10 metrics show NO full
match) are not unexpected given the severity of the experimental
intervention.  By removing the WEATHER\_DROP channel, Round~3
eliminates the mechanism that most differentiated protocol performance
in Round~2.  When all protocols achieve routing-drop rates below 2\%,
the scheduling policy's ability to handle requests \emph{that arrive}
matters less than its ability to manage dock contention and queue
pressure---a regime in which timeout rate and waiting time become the
binding constraints.  This shift explains why protocols that were
penalised in Round~2 for poor request reachability (e.g.\
\texttt{ugv\_edf}, \texttt{ugv\_p\_role\_priority}) recover
disproportionately in Round~3.

The sole partial match---energy per session ($\tau = 0.81$)---is
consistent with the expectation that per-session energy is governed by
dock hardware parameters rather than by the scheduling or weather
layer, and therefore remains invariant across rounds.

\paragraph{Transition to Chapter~5.}
Taken together, the Round~2 and Round~3 results provide
complementary views of the same system.  Round~2 reveals the full
cross-layer mechanism under realistic weather stress, identifying
which protocols are most resilient when the feedback loop is active.
Round~3 isolates the scheduling layer by removing the environmental
trigger, showing that weather variance was a necessary amplifier of
the performance gaps but not the sole source of protocol
differentiation.  Both rounds agree that survivability---measured by
battery depletions---remains the central mediator between charging
outcomes and network quality.  Chapter~5 synthesises these findings,
discusses their implications for scheduling-policy design in
shared-substrate FANET systems, and outlines directions for future
work.
